% Generated draft: requires mathematical and source review.
\documentclass[11pt]{article}
\usepackage[margin=25mm]{geometry}
\usepackage{fontspec}
\setmainfont{Latin Modern Roman}
\usepackage{amsmath,amssymb,mathtools,unicode-math}
\setmathfont{Latin Modern Math}
\usepackage{enumitem,needspace,xurl,hyperref}
\setlength{\parindent}{0pt}
\setlength{\parskip}{5pt}
\setlength{\emergencystretch}{3em}
\newcommand{\N}{\mathbb{N}}
\newcommand{\Z}{\mathbb{Z}}
\newcommand{\Q}{\mathbb{Q}}
\newcommand{\R}{\mathbb{R}}
\newcommand{\C}{\mathbb{C}}
\newcommand{\F}{\mathbb{F}}
\newcommand{\A}{\mathbb{A}}
\newcommand{\PP}{\mathbb{P}}
\newcommand{\E}{\mathbb{E}}
\newcommand{\eps}{\varepsilon}
\newcommand{\dd}{\,\mathrm{d}}
\begin{document}
\section*{UnsolvedMath ID: 5300090}
\textbf{Source problem code:} catalogue entry
\subsection{Definitions and mathematical statement}
Let \(M\) be a compact smooth \(3\)-manifold and \(f:M\to M\) a diffeomorphism. A nonempty compact invariant set \(A\) is a hyperbolic attractor when it is obtained from a trapping neighborhood and admits a continuous \(Df\)-invariant splitting
\[
T_AM=E^s\oplus E^u
\]
with constants \(C>0\), \(0<\lambda<1\), such that for all \(x\in A\) and \(n\ge0\),
\[
\|Df_x^n v\|\le C\lambda^n\|v\|\ (v\in E_x^s),\qquad
\|Df_x^{-n}w\|\le C\lambda^n\|w\|\ (w\in E_x^u).
\]
A trapping neighborhood can be expressed as an open \(U\) with \(f(\overline U)\subset U\) and \(A=\bigcap_{n\ge0}f^n(\overline U)\). Whether attractor additionally means topologically transitive is a source convention not supplied here.

The unstable lamination consists of local unstable manifolds inside \(A\). Its local product charts have an unstable plaque direction and a transverse set; a natural transversal near \(x\) is a suitable portion of \(A\cap W^s_{\mathrm{loc}}(x)\). Let \(S\) be the standard Sierpiński carpet, obtained from the unit square by repeatedly deleting open middle-third subsquares. The specific realization question is whether some such \((M,f,A)\) has a transverse structure modeled on \(S\), rather than the transverse structures in the known cases. To require an actual transversal homeomorphic to \(S\), or only local pieces of that model, one must specify the intended transversal convention.

The broader question asks whether any homeomorphism type of hyperbolic attractor occurs beyond the source's Anosov, laminated, Williams, and invariant-torus cases. Those class names do not supply a precise list of excluded homeomorphism types or their dynamical hypotheses, so an exhaustive outside-the-list predicate cannot be reconstructed from the extraction alone. No stable/unstable index, transitivity, manifold boundary condition, or additional regularity is silently imposed. The supplied research summary reports partial classification in the named families and describes additional topology, including the carpet possibility, as still under investigation.
\subsection{Short English statement}
Can a hyperbolic attractor of a compact \(3\)-manifold diffeomorphism have topology beyond the listed known classes, especially Sierpiński-carpet transversals to its unstable lamination? The class list and transversal convention need precise source definitions.
\subsection{Sources}
\begin{itemize}
\item Local UnsolvedMath registry, permanent ID \(5300090\), exact source code \texttt{catalogue entry}, ``Topology of hyperbolic attractors in dimension three.'' Proposer Christian Bonatti; proposed year 1999.
\item Stony Brook, \emph{Open Problems in Dynamical Systems}, Bonatti problem, later restatement; \url{https://www.math.stonybrook.edu/open-problems-dynamical-systems}. Recorded access 2026-07-29; source-TeX-and-pdf-text extraction.
\item Supplied research summary refers to Williams, branched/laminar, Anosov, and torus classification cases but gives no complete publication references or precise classification statements.
\end{itemize}
\end{document}
